\section{Трансдуктивный слой}
\label{sec:trans}

Про тест известно больше, чем обычно, и этим можно пользоваться. Аккуратно.

\emph{Статус раздела.} Обе описанные конструкции с тех пор упёрлись в измерение, и по-разному:
перекалибровка под тестовую маргиналь остаётся \emph{открытым решением}, которое нельзя проверить
изнутри, а серийный случайный эффект \emph{закрыт} --- ниже сказано, чем именно. Зато появилась одна
величина, которую можно посчитать \emph{на самом тесте} без единой метки, и она в конце раздела.

\subsection{Сдвиг распределения}

Тест обогащён активными по трём ферментам и близок к фону по четвёртому (раздел~\ref{sec:data}). Это ковариатный сдвиг
с индуцированным сдвигом меток: $p_{\text{тест}}(y) \ne p_{\text{обуч}}(y)$. Модель, откалиброванная
на обучающей маргинали, будет систематически занижать активность тестовых соединений. Лечится
перекалибровкой под ожидаемую тестовую маргиналь, оценённую из известной конструкции набора, и проверяется
на разбиении, воспроизводящем дизайн теста.

\subsection{Серийный случайный эффект}

Тест --- серии аналогов вокруг якорей, но охват этой структуры зависит от того, что считать серией:
при рабочем пороге 0.35 в группы по пять и более членов попадает 155 соединений из 750, при 0.48 ---
496 (раздел~\ref{sec:data}). Естественно предположить, что соединения одной серии имеют общий сдвиг:
\begin{equation}
\hat y_i = f(x_i) + b_{s(i)},
\qquad b_s \sim \mathcal{N}\bigl(\mu(\text{якорь}_s),\, \tau^2\bigr) .
\label{eq:series}
\end{equation}

Здесь $s(i)$ --- серия, к которой относится соединение $i$, $b_s$ --- её общий сдвиг, $\mu(\text{якорь}_s)$ ---
центр приора, привязанный к активности якоря, $\tau$ --- насколько сильно мы разрешаем серии отклоняться
от него.

Соблазн стянуть предсказание к активности якоря надо давить. Медианная разница $\pic$ в обучающих парах
со сходством Танимото 0.6--0.85 составляет от 0.33 до 0.68 логарифмической единицы. Обрывы активности
реальны: похожие молекулы регулярно различаются в разы. Значит $\tau$ должно быть порядка 0.5, а не 0.1,
и оценивать его надо по обучающим парам при том же уровне сходства, а не подбирать.

Оговорка, которая тут важнее самого числа: пар в этой полосе всего 601 на всю обучающую выборку, и обе
метки есть лишь у 30 из них на 1A2, 15 на 2C9, 23 на 2D6 и 510 на 3A4. То есть три из четырёх медиан
стоят на полутора-трёх десятках пар, и приводить их с двумя знаками --- преувеличение точности. Вывод
про порядок $\tau$ от этого не страдает: он одинаков при любом прочтении этих чисел. А вот оценивать
$\tau$ по обучающим парам пофермент, как здесь предлагается, на 2C9 не выйдет --- пятнадцати пар для
этого мало, и придётся либо пулить ферменты, либо расширять полосу сходства.

Проверять эффект от \eqref{eq:series} можно только на псевдотестовых сериях: на случайном разбиении
он покажет ложный выигрыш, потому что там соседи по серии лежат в обучении.

\emph{Слой закрыт, и закрыт дважды.} Первое: направление проверяется без единой тестовой метки ---
достаточно сравнить, насколько модель разводит членов серии, с тем, насколько их разводит химия.
Модель разводит \emph{слабее} в полтора-три раза на каждом ферменте, то есть она уже переглажена
внутри серии, и стягивание к якорю только усугубило бы ошибку.

Обратное лечение --- \emph{растянуть} внутрисерийные отклонения --- напрашивается и тоже не проходит,
но по другой причине, и она тоньше. Растяжение монотонно внутри группы, поэтому внутрисерийный порядок
не меняет вовсе; вся его польза межсерийная и существует ровно тогда, когда внутрисерийные отклонения
модели \emph{информативны}, а не просто сжаты. Отношение масштабов этого не различает: чистый шум дал
бы ту же картинку. Различает коэффициент затухания --- корреляция отклонений, умноженная на отношение
размахов, --- и на CYP3A4, единственном ферменте, где серий хватает для оценки, он равен 0.833. Ниже
единицы: отклонения зашумлены сильнее, чем сжаты, и растяжение усилило бы шум.

Второе, и это ограничение на весь раздел: \emph{наша кросс-валидация к серийной оси структурно слепа}.
При сходстве 0.60 в сериях лежит 10.2 процента обучающих молекул против 69.1 процента тестовых ---
всемеро. Любая оценка серийного эффекта, снятая на нашей выборке, ограничивает не тест, а популяцию, в
которой серий почти нет.

\subsection{Что о тесте всё же можно узнать без меток}

Одна величина считается прямо на 750 закрытых молекулах и меток не требует: доля пар, о порядке которых
члены ансамбля \emph{расходятся}. Спорна пара или нет --- вопрос о том, одинаково ли её упорядочили
пять моделей, а не о том, кто прав.

Вне фолда таких пар 27.4 процента, на тесте 27.9 --- отношение 1.012. Подача, стало быть, ровно
настолько же внутренне согласована на тесте, насколько была на кросс-валидации, и точность на спорных
парах переносится как измерена. Это отрицательный результат про опасение, а не выигрыш; но после того
как выяснилось, что к аналоговой структуре теста наша проверка слепа, утверждение о тесте, не
требующее меток, стоит здесь дороже обычного.
